\documentclass{article}
\usepackage[T1]{fontenc}
\usepackage{lmodern}
\usepackage{graphicx}
\usepackage{amsmath,amssymb}
\usepackage[a4paper,margin=2cm]{geometry}
\usepackage[absolute,overlay]{textpos}
\setlength{\TPHorizModule}{1mm}
\setlength{\TPVertModule}{1mm}
\usepackage[indonesian]{babel}
\usepackage{hyperref}
\setlength{\emergencystretch}{2em}
% unit: Halaman 4

\begin{document}

% === Halaman 4 (llm) ===

\begin{enumerate}
  \item[3)] Alice membangkitkan kunci rahasia (K) untuk enkripsi pesan dengan AES
  \item[4)] Alice mengenkripsi K dengan RSA menggunakan kunci publik Bob (PK$_{Bob}$).
  \[ E\_RSA_{PKbob}(K) = CK \]
  \item[5)] Alice mengenkripsi pesan M dengan AES menggunakan K,
  \[ E\_AES_{K}(M) = CM \]
  lalu mengirim CK dan CM kepada Bob.
  \item[6)] Bob mendekripsi CK dengan RSA menggunakan kunci privatnya (SK$_{Bob}$)
  \[ D\_RSA_{SKbob}(CK) = K \]
  \item[7)] Selanjutnya Bob mendekripsi pesan CM dengan AES menggunakan K
  \[ D\_AES_{K}(CM) = M \]
\end{enumerate}

\vspace{10mm}
\centering{Rinaldi M/II4020 Kriptografi}

\end{document}
